\begin{frame}{어텐션의 직관: 모든 단어를 한 번에}
  ``지금 이 단어를 이해하려면, 문장의 \textbf{다른 모든 단어를 한
  번에} 보고, 그중 관련 있는 단어에 더 집중(attend)한다.''
  \vskip0.6em
  예: \textbf{``그 동물은 도로를 건너지 않았다. 왜냐하면 그것은
  너무 지쳐 있었기 때문이다.''} — ``그것은''이 ``그 동물''인지
  ``도로''인지는 \textbf{문장 전체를 동시에 봐야} 판단할 수 있다.
  \vskip0.5em
  \begin{block}{RNN과의 차원이 달라지는 수식}
    어텐션은 현재 단어의 새 표현을 만들 때 $n$개 단어 각각에
    \textbf{가중치 하나씩, 총 $n$개}를 계산 — ``단어 위의 분포''
    를 만든다. 이 $n$차원 벡터를 \kb{어텐션 가중치}라 부른다.
    ``그것은''의 가중치는 (0.56, 0.01, 0.01, 0.42) — ``동물''에
    56\%, 자기 자신에 42\% (숫자는 12.2절 실습).
    \vskip0.3em
    RNN이 주던 것: ``과거의 \textbf{압축 요약(한 벡터)}''.
    어텐션이 주는 것: \textbf{``관련 단어들의 내용을 적절한 비율로
    섞은 것''}.
  \end{block}
  \vskip0.4em
  \kq{어텐션의 기원은 새로운 아키텍처가 아니라, RNN에 붙인
  ``패치''였다} — 3년 뒤 트랜스포머는 그 메커니즘 하나만 남기고
  RNN 본체를 통째로 걷어냈다. 이 장은 Week 11 서사의 귀결이다.
\end{frame}
